%%%%%%%%%%%%%%%%%%%%%%%%%%%%%%%%%%%%%%%%%%%%%%%%%%%%%%%%%%%%%%%%%%%%%%
% Lớp 8 | Chương II. Hằng đẳng thức đáng nhớ và ứng dụng | Hiệu hai bình phương. Bình phương của một tổng hay một hiệu
% Dạng 8D2?6-1: Khai triển hằng đẳng thức
%%%%%%%%%%%%%%%%%%%%%%%%%%%%%%%%%%%%%%%%%%%%%%%%%%%%%%%%%%%%%%%%%%%%%%


% nguồn: dan_tu_AI (nhập 14/11/2026)
\begin{bt}%[8D2H6-1]%[ma:9c995729]
	Khai triển các biểu thức sau:\\
	a) $(x+3)^2$;\qquad b) $(2y-5)^2$;\qquad c) $(3a-b)(3a+b)$.
	\loigiai{
		a) $(x+3)^2=x^2+2\cdot x\cdot 3+3^2=x^2+6x+9$.\\
		b) $(2y-5)^2=(2y)^2-2\cdot 2y\cdot 5+5^2=4y^2-20y+25$.\\
		c) $(3a-b)(3a+b)=(3a)^2-b^2=9a^2-b^2$.
	}
\end{bt}

% nguồn: dan_tu_AI (nhập 14/11/2026)
\begin{bt}%[8D2H6-1]%[ma:9340f162]
	Khai triển các biểu thức sau:\\
	a) $\left(\dfrac12x+2\right)^2$;\qquad b) $\left(x^2-3y\right)^2$;\qquad c) $(5-4x)(4x+5)$.
	\loigiai{
		a) $\left(\dfrac12x+2\right)^2=\left(\dfrac12x\right)^2+2\cdot\dfrac12x\cdot 2+2^2=\dfrac14x^2+2x+4$.\\
		b) $\left(x^2-3y\right)^2=\left(x^2\right)^2-2\cdot x^2\cdot 3y+(3y)^2=x^4-6x^2y+9y^2$.\\
		c) $(5-4x)(4x+5)=(5-4x)(5+4x)=5^2-(4x)^2=25-16x^2$.
	}
\end{bt}

% nguồn: dan_tu_AI (nhập 14/11/2026)
\begin{bt}%[8D2H6-1]%[ma:dabbc849]
	Rút gọn các biểu thức sau:\\
	a) $A=(x+2)^2+(x-2)^2$;\\
	b) $B=(x+2)^2-(x-2)^2$;\\
	c) $C=(x+1)(x-1)-(x-1)^2$.
	\loigiai{
		a) $A=\left(x^2+4x+4\right)+\left(x^2-4x+4\right)=2x^2+8$.\\
		b) $B=\left(x^2+4x+4\right)-\left(x^2-4x+4\right)=x^2+4x+4-x^2+4x-4=8x$.\\
		c) $C=\left(x^2-1\right)-\left(x^2-2x+1\right)=x^2-1-x^2+2x-1=2x-2$.
	}
\end{bt}

% nguồn: dan_tu_AI (nhập 14/11/2026)
\begin{bt}%[8D2V6-1]%[ma:54a58c79]
	a) Bằng cách viết $a+b+c=(a+b)+c$, hãy chứng tỏ
	\[ (a+b+c)^2=a^2+b^2+c^2+2ab+2bc+2ca. \]
	b) Áp dụng, khai triển biểu thức $(x+2y-3)^2$.
	\loigiai{
		a) Áp dụng bình phương của một tổng hai lần:
		\[ (a+b+c)^2=\left[(a+b)+c\right]^2=(a+b)^2+2(a+b)c+c^2=a^2+2ab+b^2+2ac+2bc+c^2. \]
		Vậy $(a+b+c)^2=a^2+b^2+c^2+2ab+2bc+2ca$.\\
		b) Áp dụng với $a=x$, $b=2y$, $c=-3$:
		\[ (x+2y-3)^2=x^2+(2y)^2+(-3)^2+2\cdot x\cdot 2y+2\cdot 2y\cdot(-3)+2\cdot(-3)\cdot x \]
		\[ =x^2+4y^2+9+4xy-12y-6x. \]
	}
\end{bt}

% nguồn: dan_tu_AI (nhập 14/11/2026)
\begin{bt}%[8D2V6-1]%[ma:c3ee6708]
	Một mảnh vườn hình vuông có cạnh $x$ mét ($x>3$).\\
	a) Nếu mỗi cạnh tăng thêm $3\text{ m}$ thì được mảnh vườn hình vuông mới. Viết biểu thức (dạng khai triển) biểu thị diện tích mảnh vườn mới. Diện tích tăng thêm bao nhiêu mét vuông?\\
	b) Nếu một cạnh tăng thêm $3\text{ m}$, cạnh kia giảm đi $3\text{ m}$ thì được mảnh vườn hình chữ nhật. Diện tích mảnh vườn hình chữ nhật này lớn hơn hay nhỏ hơn diện tích mảnh vườn ban đầu bao nhiêu mét vuông?
	\loigiai{
		a) Diện tích mảnh vườn mới:
		\[ (x+3)^2=x^2+6x+9\ \left(\text{m}^2\right). \]
		Diện tích tăng thêm: $\left(x^2+6x+9\right)-x^2=6x+9\ \left(\text{m}^2\right)$.\\
		b) Diện tích mảnh vườn hình chữ nhật:
		\[ (x+3)(x-3)=x^2-9\ \left(\text{m}^2\right). \]
		Vậy diện tích mảnh vườn hình chữ nhật nhỏ hơn diện tích mảnh vườn ban đầu $9\text{ m}^2$.
	}
\end{bt}

% nguồn: dan_tu_AI (nhập 16/11/2026)
\begin{bt}%[8D2V6-1]%[ma:36793ee8]
	Bác An có một mảnh vườn hình vuông cạnh $x+4$ (m). Bác đào một cái ao hình vuông cạnh $x-2$ (m) ở giữa vườn ($x>2$). Viết biểu thức tính diện tích phần đất còn lại và khai triển biểu thức đó.
	\loigiai{
		Diện tích mảnh vườn hình vuông là: $S_1 = (x+4)^2$ ($\text{m}^2$).\\
		Diện tích cái ao hình vuông là: $S_2 = (x-2)^2$ ($\text{m}^2$).\\
		Biểu thức tính diện tích phần đất còn lại là:
		\[ S = S_1 - S_2 = (x+4)^2 - (x-2)^2. \]
		Khai triển biểu thức $S$:
		\begin{align*} 
			S &= (x^2 + 8x + 16) - (x^2 - 4x + 4) \\
			&= x^2 + 8x + 16 - x^2 + 4x - 4 \\
			&= 12x + 12 \text{ (m}^2\text{)}.
		\end{align*}
		(Hoặc dùng hiệu hai bình phương: $S = [(x+4)-(x-2)][(x+4)+(x-2)] = 6(2x+2) = 12x+12$).
	}
\end{bt}

% nguồn: De_AI_Toan8_cuoi_hoc_ki_i_1 (nhập 21/11/2026)
\begin{bt}%[8D2N6-1]%[ma:39d59352]
	Khai triển các hằng đẳng thức sau:
a) $(x + 3)^2$;
b) $(2x - 1)^2$.
	\loigiai{
		a) Áp dụng hằng đẳng thức bình phương của một tổng $(A + B)^2 = A^2 + 2AB + B^2$ với $A = x$, $B = 3$: \\
$(x + 3)^2 = x^2 + 2 \cdot x \cdot 3 + 3^2 = x^2 + 6x + 9$. \\
b) Áp dụng hằng đẳng thức bình phương của một hiệu $(A - B)^2 = A^2 - 2AB + B^2$ với $A = 2x$, $B = 1$: \\
$(2x - 1)^2 = (2x)^2 - 2 \cdot 2x \cdot 1 + 1^2 = 4x^2 - 4x + 1$.
	}
\end{bt}

% nguồn: De_AI_Toan8_cuoi_hoc_ki_i_1 (nhập 21/11/2026)
\begin{bt}%[8D2N6-1]%[ma:07035abc]
	Khai triển các biểu thức sau bằng hằng đẳng thức:
a) $(x - 5)(x + 5)$;
b) $(3x + 2y)^2$.
	\loigiai{
		a) Áp dụng hằng đẳng thức hiệu hai bình phương $(A - B)(A + B) = A^2 - B^2$ với $A = x$, $B = 5$: \\
$(x - 5)(x + 5) = x^2 - 5^2 = x^2 - 25$. \\
b) Áp dụng hằng đẳng thức bình phương của một tổng với $A = 3x$, $B = 2y$: \\
$(3x + 2y)^2 = (3x)^2 + 2 \cdot 3x \cdot 2y + (2y)^2 = 9x^2 + 12xy + 4y^2$.
	}
\end{bt}

% nguồn: De_AI_Toan8_cuoi_hoc_ki_i_1 (nhập 21/11/2026)
\begin{bt}%[8D2N6-1]%[ma:2e3f941a]
	Khai triển các biểu thức sau:
a) $(4x - 3y)^2$;
b) $(2x - 5y)(2x + 5y)$.
	\loigiai{
		a) Áp dụng hằng đẳng thức bình phương của một hiệu với $A = 4x$, $B = 3y$: \\
$(4x - 3y)^2 = (4x)^2 - 2 \cdot 4x \cdot 3y + (3y)^2 = 16x^2 - 24xy + 9y^2$. \\
b) Áp dụng hằng đẳng thức hiệu hai bình phương với $A = 2x$, $B = 5y$: \\
$(2x - 5y)(2x + 5y) = (2x)^2 - (5y)^2 = 4x^2 - 25y^2$.
	}
\end{bt}

% nguồn: De_AI_Toan8_cuoi_hoc_ki_i_1 (nhập 21/11/2026)
\begin{bt}%[8D2N6-1]%[ma:0656498e]
	Khai triển các hằng đẳng thức sau:
a) $(x + 2)^3$;
b) $(x - 3)^3$.
	\loigiai{
		a) Áp dụng hằng đẳng thức lập phương của một tổng $(A + B)^3 = A^3 + 3A^2B + 3AB^2 + B^3$: \\
$(x + 2)^3 = x^3 + 3 \cdot x^2 \cdot 2 + 3 \cdot x \cdot 2^2 + 2^3 = x^3 + 6x^2 + 12x + 8$. \\
b) Áp dụng hằng đẳng thức lập phương của một hiệu $(A - B)^3 = A^3 - 3A^2B + 3AB^2 - B^3$: \\
$(x - 3)^3 = x^3 - 3 \cdot x^2 \cdot 3 + 3 \cdot x \cdot 3^2 - 3^3 = x^3 - 9x^2 + 27x - 27$.
	}
\end{bt}

% nguồn: De_AI_Toan8_cuoi_hoc_ki_i_1 (nhập 21/11/2026)
\begin{bt}%[8D2N6-1]%[ma:1be212a7]
	Khai triển các biểu thức sau:
a) $(2x + 1)^3$;
b) $(x - 2y)^3$.
	\loigiai{
		a) Áp dụng hằng đẳng thức lập phương của một tổng với $A = 2x$, $B = 1$: \\
$(2x + 1)^3 = (2x)^3 + 3 \cdot (2x)^2 \cdot 1 + 3 \cdot 2x \cdot 1^2 + 1^3 = 8x^3 + 12x^2 + 6x + 1$. \\
b) Áp dụng hằng đẳng thức lập phương của một hiệu với $A = x$, $B = 2y$: \\
$(x - 2y)^3 = x^3 - 3 \cdot x^2 \cdot 2y + 3 \cdot x \cdot (2y)^2 - (2y)^3 = x^3 - 6x^2y + 12xy^2 - 8y^3$.
	}
\end{bt}

% nguồn: De_AI_Toan8_cuoi_hoc_ki_i_1 (nhập 21/11/2026)
\begin{bt}%[8D2N6-1]%[ma:e7c1757c]
	Khai triển các biểu thức sau:
a) $\left(x + \dfrac{1}{2}\right)^2$;
b) $\left(x - \dfrac{1}{3}\right)\left(x + \dfrac{1}{3}\right)$.
	\loigiai{
		a) Áp dụng hằng đẳng thức bình phương của một tổng: \\
$\left(x + \dfrac{1}{2}\right)^2 = x^2 + 2 \cdot x \cdot \dfrac{1}{2} + \left(\dfrac{1}{2}\right)^2 = x^2 + x + \dfrac{1}{4}$. \\
b) Áp dụng hằng đẳng thức hiệu hai bình phương: \\
$\left(x - \dfrac{1}{3}\right)\left(x + \dfrac{1}{3}\right) = x^2 - \left(\dfrac{1}{3}\right)^2 = x^2 - \dfrac{1}{9}$.
	}
\end{bt}

% nguồn: De_AI_Toan8_cuoi_hoc_ki_i_1 (nhập 21/11/2026)
\begin{bt}%[8D2N6-1]%[ma:34d6616f]
	Thu gọn các biểu thức sau:
a) $A = (x + 2)^2 - x^2$;
b) $B = (x - 3)^2 - (x^2 + 9)$.
	\loigiai{
		a) $A = (x + 2)^2 - x^2 = (x^2 + 4x + 4) - x^2 = 4x + 4$. \\
b) $B = (x - 3)^2 - (x^2 + 9) = (x^2 - 6x + 9) - x^2 - 9 = -6x$.
	}
\end{bt}

% nguồn: De_AI_Toan8_cuoi_hoc_ki_i_1 (nhập 21/11/2026)
\begin{bt}%[8D2N6-1]%[ma:1c30fcc5]
	Thu gọn các biểu thức sau:
a) $C = (x + 4)(x - 4) - x^2$;
b) $D = (2x + 1)^2 + (2x - 1)^2$.
	\loigiai{
		a) $C = (x + 4)(x - 4) - x^2 = (x^2 - 16) - x^2 = -16$. \\
b) $D = (2x + 1)^2 + (2x - 1)^2 = (4x^2 + 4x + 1) + (4x^2 - 4x + 1) = 8x^2 + 2$.
	}
\end{bt}

% nguồn: De_AI_Toan8_cuoi_hoc_ki_i_1 (nhập 21/11/2026)
\begin{bt}%[8D2N6-1]%[ma:19df25d4]
	Thu gọn các biểu thức sau:
a) $E = (x + 3)^2 - (x - 3)^2$;
b) $F = (x + 1)^2 - (x - 2)(x + 2)$.
	\loigiai{
		a) $E = (x + 3)^2 - (x - 3)^2 = (x^2 + 6x + 9) - (x^2 - 6x + 9) = x^2 + 6x + 9 - x^2 + 6x - 9 = 12x$. \\
b) $F = (x + 1)^2 - (x - 2)(x + 2) = (x^2 + 2x + 1) - (x^2 - 4) = x^2 + 2x + 1 - x^2 + 4 = 2x + 5$.
	}
\end{bt}
